\documentclass[11pt]{amsart}
\usepackage[T1]{fontenc}
\usepackage[utf8]{inputenc}
\usepackage{amssymb}
\usepackage[margin=1.1in]{geometry}
\usepackage[hidelinks]{hyperref}
\newcommand{\stamptime}{\the\numexpr\time/60\relax:\ifnum\numexpr\time-60*(\time/60)\relax<10 0\fi\the\numexpr\time-60*(\time/60)\relax}
\AddToHook{shipout/foreground}{\put(0.5\paperwidth,-\paperheight+0.55cm){\makebox[0pt]{\tiny\ttfamily unit/doc/\jobname.tex \ --- \ compiled \today, \stamptime}}}
\newtheorem{theorem}{Theorem}
\newtheorem{proposition}[theorem]{Proposition}
\theoremstyle{remark}
\newtheorem{remark}[theorem]{Remark}
\newcommand{\code}[1]{\texttt{#1}}
\newcommand{\Om}{\Omega}
\newcommand{\Fib}{\mathrm{Fib}}
\newcommand{\Map}{\mathrm{Map}_*}
\newcommand{\Mor}{\mathsf{Mor}}
\newcommand{\Hom}{\mathrm{Hom}}
\newcommand{\Tri}{\mathsf{Triples}}
\newcommand{\SCt}{\mathsf{SC3}}
\newcommand{\SCf}{\mathsf{SC4}}
\newcommand{\Z}{\mathbb{Z}}
\newcommand{\cmp}{\mathbin{;}}
\newcommand{\real}{\mathsf{realise}}
\title{Postnikov triples, their maps, and the spheres:\\ an overview of the argument}
\date{28 October 2026}
\begin{document}
\maketitle

\noindent This note describes, at a high level, the argument formalised in
\code{sag/pi3-s2-postnikov} by the modules \code{TripleMaps},
\code{TripleMapsEquiv}, \code{TripleComponents}, \code{TripleHomotopy},
\code{Pi3TripleUniversal} and \code{Pi4TripleUniversal}. The detailed informal
proof, with a declaration-by-declaration correspondence, is
\code{doc/postnikov-morphisms.pdf}.

\noindent\textbf{Attribution.} The strategy is David Wärn's, as communicated
by the user: classify pointed simply connected $3$-types \emph{and their
arbitrary pointed maps} by Postnikov data, identify the universal objects of the
spheres on the side of the data, and treat $\pi_4(S^3)$ by embedding pointed
$2$-connected $4$-types into pointed simply connected $3$-types by looping. The
classification of objects is Hart--Milner's (Owen Milner's equivalence). We do
not know Owen Milner's unpublished treatment of maps; nothing here claims to
reproduce it.

\section{The plan}

To compute $\pi_3(S^2)$ it suffices to know a pointed simply connected $3$-type
$P$ with a class $s\in[S^2,P]_*$ such that
\[s^*:[P,X]_*\to[S^2,X]_*\]
is a bijection for every pointed simply connected $3$-type $X$: then $P$ is the
$3$-truncation of $S^2$, and $\pi_3(S^2)=\pi_3(P)$ by a Yoneda argument. Similarly,
$\pi_4(S^3)$ is computed from a pointed $2$-connected $4$-type with the analogous
universal property among such types.

Wärn's idea is to make the left-hand side of this bijection \emph{algebraic}.
If pointed simply connected $3$-types and their maps are described by algebraic
data, then $[P,X]_*$ becomes a set of classes of algebraic morphisms, computable
from the cohomology of Eilenberg--MacLane spaces, while $[S^2,X]_*$ is computed
from the sphere. The argument has three steps:
\begin{enumerate}
\item classify the objects (done earlier: \code{PostnikovClassification});
\item classify the arbitrary pointed maps (new: Theorem~\ref{thm:maps});
\item compute the sphere's universal property on the algebraic side (new:
Sections~\ref{sec:pi3} and~\ref{sec:pi4}), and recover the homotopy group.
\end{enumerate}

\section{Objects}

A \emph{Postnikov triple} is $t=(A,B,k)$ with abelian groups $A,B$ and a pointed
map $k:K(A,2)\to_*K(B,4)$. Its realisation is the fibre
\[\Fib k=\textstyle\sum_{x:K(A,2)}(k\,x=\mathrm{pt}),\]
a pointed simply connected $3$-type with $\pi_2=A$ and $\pi_3=B$. The map
\[\real:\Tri\to\SCt\]
is an equivalence of types (\code{PostnikovClassification.classification}).
Both types live in the next universe and are not sets; they are $2$-types.

\section{Maps}\label{sec:maps}

For triples $t=(A,B,k)$ and $t'=(A',B',k')$ define
\[\Mor(t,t')=\sum_{a:\Hom(A,A')}\ \sum_{b:\Hom(B,B')}\big(K(a,2)\cmp k'=k\cmp K(b,4)\big).\]
Here $f\cmp g$ means ``$f$, then $g$'', and the last component is a
\emph{path of pointed maps}: a homotopy together with its coherence at the base
point. Such a morphism determines a map of fibres, $(x,q)\mapsto(K(a,2)\,x,\ H_x\cdot
K(b,4)(q))$. This is the canonical construction \code{FibMap.fm} already present
in the library.

\begin{theorem}[\code{TripleMapsEquiv.Maps2.realise-map-equiv}]\label{thm:maps}
For all triples $t,t'$, the canonical map
\[\Mor(t,t')\longrightarrow\Map(\Fib k,\Fib k')\]
is an equivalence of types. There is no hypothesis on the groups or on $k,k'$,
and no sphere axiom.
\end{theorem}

Both sides are $1$-types, not sets in general. The homotopy $H$ is genuine data.
For $t=(A,0,0)$ and $t'=(0,C,0)$ the coefficients carry no information at all,
and both sides are $\Map(K(A,2),K(C,3))$, whose components form
$H^3(K(A,2);C)$.

\subsection*{Idea of the proof}
A pointed map $\Fib k\to\Fib k'$ is a map to $K(A',2)$ together with a pointed
null-homotopy of its composite with $k'$.
\begin{itemize}
\item \emph{The base.} The projection $\Fib k\to K(A,2)$ has $2$-connected
fibres, and $K(A',2)$ is a $2$-type. So maps into $K(A',2)$ factor uniquely
through $K(A,2)$, and $\Map(K(A,2),K(A',2))\simeq\Hom(A,A')$.
\item \emph{The key lemma.} For $c:K(A,2)\to_*K(B',4)$, pointed null-homotopies
of $\Fib k\to K(A,2)\xrightarrow{c}K(B',4)$ are equivalent to factorisations of
$c$ through $k$, that is, to pairs $(\beta,\ c=k\cmp\beta)$
(\code{TripleMaps.Key.key}).
\end{itemize}
In fibrewise terms, a null-homotopy is a family of maps
$F_x:(k\,x=\mathrm{pt})\to(c\,x=\mathrm{pt})$ between ``$K(B,3)$-torsors''. The
proof separates a linear part from a translation:
\begin{itemize}
\item whether $F_x$ is a translate of the map induced by a fixed $\beta$ is a
\emph{proposition}, by the connectivity of the fibres and the truncation of
the target;
\item so it suffices to check it over the base point, since the base is
connected;
\item over the base point, the restriction of $F$ is a pointed map
$\Om K(B,4)\to\Om K(B',4)$, and it has a \emph{unique} delooping $\beta$
(Wärn's delooping theorem, \code{DeloopMaps}, Corollary~11);
\item the translations are exactly the homotopy component, which stays as data.
\end{itemize}
No point is chosen in any fibre other than the base, no representative of a
class is chosen, and lifts are not assumed unique. Finally,
$\Map(K(B,4),K(B',4))\simeq\Hom(B,B')$. The composite of these equivalences is,
by definition, the canonical map \code{FibMap.fm}.

\section{The homotopy precategory of triples}\label{sec:cat}

Composition of morphisms of triples is defined intrinsically. The coefficients
compose, and the homotopies paste:
\[K(a\cmp a',2)\cmp k''=K(a,2)\cmp k'\cmp K(b',4)=k\cmp K(b\cmp b',4).\]
The identity uses $K(1,m)=\mathrm{id}$. The canonical maps of fibres preserve
composition and identities up to pointed homotopy, that is, in the homotopy
classes $[\Fib k,\Fib k'']_*$ (\code{TripleHomotopy.Maps3.rmap-comp},
\code{Maps1.rmap-id}). Put $\hom(t,t')=\|\Mor(t,t')\|_0$.

\begin{theorem}[\code{TripleHomotopy.Cat}, with \code{PostnikovClassification}]\label{thm:hocat}
Triples, with the sets $\hom(t,t')$ and the induced composition and identities,
form a precategory. The functor to the homotopy precategory of pointed simply
connected $3$-types has two parts:
\begin{itemize}
\item on objects it is $\real$, an equivalence of types;
\item on hom-sets it is the class of the canonical map, a bijection by
Theorem~\ref{thm:maps}.
\end{itemize}
It preserves composition and identities.
\end{theorem}

In the terminology of Ahrens, Kapulkin and Shulman
(\emph{Univalent categories and the Rezk completion}, Definition~6.9), this is
an \emph{isomorphism of precategories}: a fully faithful functor whose object map
is an equivalence of types. A precategory may have an arbitrary type of objects
(here $2$-types), while its hom-types are sets. Neither precategory is univalent
in their sense, and no such claim is made. By their Lemma~6.16, under univalence
an isomorphism of precategories yields an equality of the two precategories;
this consequence is not formalised here. No explicitly declared inverse functor
is needed for the criterion. Equality in $[\Fib k,\Fib k'']_*$ is all the
homotopy precategory requires, so no coherent pointed path for the preservation
of composition is constructed.

\section{$\pi_3(S^2)$ on the triple side}\label{sec:pi3}

Take $t_0=(\Z,\Z,\kappa)$, with $\kappa$ the cup square, and $P=\Fib\kappa$.

\begin{proposition}[\code{Pi3TripleUniversal}]
For every triple $t'=(A',B',k')$, the map $\hom(t_0,t')\to\Hom(\Z,A')$,
$[a,b,H]\mapsto a$, is a bijection.
\end{proposition}

The proof is a direct computation on $\Mor$:
\begin{itemize}
\item (H4), the bilinear classification of $[K(\Z,2),K(B',4)]_*$, shows that
for each $a$ there is exactly one $b$ for which $K(a,2)\cmp k'$ and
$\kappa\cmp K(b,4)$ are homotopic;
\item (H3), $[K(\Z,2),K(B',3)]_*=0$, shows that the homotopies between them form
a single component.
\end{itemize}
The second point needs loops at an arbitrary point of
$\Map(K(\Z,2),K(B',4))$. Writing $K(B',4)\simeq\Om K(B',5)$, pointwise
multiplication by $g^{-1}$ translates them to loops at $0$
(\code{TripleComponents}).

On the sphere side, maps $S^2\to\Fib k'$ correspond to maps $S^2\to K(A',2)$,
hence to $\Hom(\Z,A')$. The marked class $s:[S^2,P]_*$ makes the triangle
commute, because the canonical map of fibres lies over $K(a,2)$. Therefore
$s^*:[P,\Fib k']_*\to[S^2,\Fib k']_*$ is a bijection for every triple. By the
classification of objects, this is the universal property for every pointed
simply connected $3$-type, and $\pi_3(S^2)\cong\Z$ follows
(\code{Pi3TripleUniversal.Final.pi3-S2}, under the $S^2$ axiom only).

\section{$\pi_4(S^3)$ on the triple side}\label{sec:pi4}

\emph{Looping} sends a pointed $2$-connected $4$-type $Y$ to the pointed simply
connected $3$-type $\Om Y$; it is an embedding $\SCf\hookrightarrow\SCt$. Its
image consists of the \emph{deloopable} triples. On maps, looping is a bijection
$[X,Y]_*\simeq[\Om X,\Om Y]_*$. The candidate is $\bar t=(\Z,\Z/2,\bar k)$, with
$\bar k$ the cup square reduced mod~$2$, and $R=\Fib\bar k=\Om BR$.

\begin{proposition}[\code{Pi4TripleUniversal}]
For every deloopable triple $t'$, the map $\hom(\bar t,t')\to\Hom(\Z,A')$ is a
bijection.
\end{proposition}

Again the computation is on $\Mor$:
\begin{itemize}
\item (H4) and Statement~T (deloopability forces $2h(1)=0$) show that the
required class factors through $\Z/2$;
\item injectivity of $g\mapsto[\bar k\cmp K(g,4)]$ fixes $b$;
\item (H3) controls the homotopies, as in Section~\ref{sec:pi3}.
\end{itemize}
On the sphere side, the source is $\Om S^3$. Lifting there needs an existence
hypothesis, and it is obtained from the surjectivity in the proposition above,
so it does not use the earlier specialised lifting lemma at $R$. The looped
marked class then has the universal property against every $\Om Y$, and by the
looping embedding statement M holds for every pointed $2$-connected $4$-type.
Both instances used in the final step come from it: $BR$ itself and
$\|S^3\|_4$. This gives $\pi_4(S^3)\cong\pi_4(BR)\cong\Z/2$
(\code{Pi4TripleUniversal.Final4.pi4-S3}, under the $S^3$ axiom only).

\section{What is reused, and what is independent}

The new sphere proofs reuse substantial mathematics from the earlier
development:
\begin{itemize}
\item the cup square and its normalisation;
\item (H3) and (H4);
\item the cross-term and two-torsion analysis (Statement~T);
\item the delooping embedding and the looped form of statement~M;
\item the sphere-side computations at $S^2$, $S^3$ and $\Om S^3$, together with
the marked classes.
\end{itemize}
They do \emph{not} use the earlier specialised mapping theorems as the reason the
universal property holds. This was checked in a scratch copy. There the old
results were deleted: \code{MapFib.Final.rep-fib}, the lifting lemma at $P$, the
lifting at $R$ and $\Om S^3$ through \code{MarkDirect.ExR}, \code{MarkDirect}'s
marked theorem, \code{DeloopEmbedding.Mark.marked}, \code{FibreH3.fib-h3}, the
presentation-based endpoints, and \code{Pi3Classified}/\code{Pi4Direct}. The copy
still checks with all the new modules. The earlier endpoints remain unchanged in
the development, as references.

\section{Status and remaining qualifications}

Everything above is checked. The whole development checks in a fresh process:
5633 declarations, 0 failed. All 32 \code{.opens} specifications pass the uniform
checker. Two qualifications remain. Neither is a missing step in the checked
sphere arguments.
\begin{itemize}
\item The theorem shows that the canonical map lies over $K(a,2)$, and that it
restricts on the base fibre to $\Om K(b,4)$. That the coefficient $b$ is the map
induced on $\pi_3$, through the identification $\pi_3(\Fib k)\cong B$, is not
formalised; it would make the interpretation of the classification more
explicit.
\item The six new \code{.opens} specifications were generated from the checker
and audited against the informal proof \emph{by kind}. Passing generated
specifications is not the entry-by-entry unfolding audit, which remains to be
done.
\end{itemize}

\end{document}
